\subsection{Low-Dimensional State and Event-Locality Tests}

After T1 survival was established as a candidate residual, we tested whether it
could be reduced to a small set of state variables. Let $C(t)$ denote the
compact-density coordinate, $\dot{C}(t)$ its temporal derivative, and $R(t)$ a
swarm-size coordinate. A low-dimensional moment-closure test asked whether
conditioning on $(C,\dot{C},R)$ improved first- or second-moment prediction of
the T1-related residual under grouped out-of-sample validation. The baseline
was a simpler radius-only model, and shifted-state controls were used to check
whether any apparent improvement exceeded a temporal-null reference.

We also tested whether true transition timestamps carried event-local
information beyond continuous state. For this analysis, T1 activity in the
near-pre window around true transition times was compared with
same-observation non-event frames matched in $(C,\dot{C},R)$. A positive result
would require true events to show robust excess activity beyond those matched
controls.

\subsection{Recent-History and Heterogeneity Tests}

Finally, we asked whether recent state-path direction supplied information
not contained in the current state. Frames were matched to have similar current
state but different recent path directions. The resulting history-conditioned
separation was compared with within-observation shuffled-history controls.
Because the sign and order of the history effect varied across observations,
this analysis was treated as an observation-specific boundary rather than a
universal memory test.

All cross-observation summaries retained the 19-observation structure. Small-n
metadata associations were used descriptively only. In particular, recording
condition and observation order were not used as causal explanations for
failure or survivor classes.

